%%%%%%%%%%  scan idx 450  |  folio 439  %%%%%%%%%%

\medskip
\noindent
10. \underline{Propiedades del acoplamiento de monodromía}.
\underline{Teorema de integralidad y positividad}

En esta sección daremos primero ciertas propiedades formales,
esencialmente triviales, de los acoplamientos de monodromía, dejando al
lector los detalles, algo tediosos, de las verificaciones. Enunciaremos
también un teorema más profundo (10.4) sobre estos acoplamientos e
indicaremos cómo reducir su demostración a un caso particular de
jacobianas, que será tratado mediante métodos trascendentes en la
sección 12, lógicamente independiente de la sección 11.

\medskip
\noindent
10.1. \underline{Transporte de estructura}. De la definición del
acoplamiento de monodromía (9.1.2) se desprende claramente que este
acoplamiento es compatible con el transporte de estructura en la
situación $(S,A_K,A'_K)$.

\medskip
\noindent
10.2. \underline{Simetría}. Partiendo del esquema dual $A'_K$ de
$A_K$ y aplicándole la construcción de la sección 9, hay que introducir
el esquema dual $(A'_K)'$ de $A'_K$, canónicamente isomorfo a $A_K$
por el teorema de bidualidad ([5] o VIII 3.2). Denotemos por
$u_\ell^{A_K}$ el acoplamiento de monodromía (9.1.2) relativo a $A_K$.
Del mismo modo obtenemos un acoplamiento
\begin{equation}\tag{10.2.1}
u_\ell^{A'_K}:\underline M'_\ell\otimes\underline M_\ell
 \longrightarrow\mathbb Z_\ell.
\end{equation}
Afirmo que este último se obtiene de $u_\ell^{A_K}$ mediante la
simetría
\begin{equation*}
s:\underline M'_\ell\otimes\underline M_\ell
 \xrightarrow{\ \sim\ }
 \underline M_\ell\otimes\underline M'_\ell;
\end{equation*}
es decir,

%%%%%%%%%%  scan idx 451  |  folio 440  %%%%%%%%%%

\begin{equation}\tag{10.2.2}
u_\ell^{A'_K}=u_\ell^{A_K}\circ s.
\end{equation}

Teniendo en cuenta la propiedad conocida de anticonmutatividad de los
acoplamientos $\varphi_\xi$ ([5] o VIII 2.2.11), la fórmula de simetría
(10.2.2) puede reformularse como sigue en términos de una polarización
fija $\xi$ de $A_K$. Esta define un homomorfismo
\begin{equation*}
\widetilde\xi:A_K\longrightarrow A'_K
\end{equation*}
y, por tanto, un homomorfismo correspondiente ($t$ significa
``tórico'')
\begin{equation}\tag{10.2.3}
\xi_t:\underline M\longrightarrow\underline M'.
\end{equation}
Podemos formar entonces el acoplamiento
\begin{equation}\tag{10.2.4}
\begin{aligned}
\varphi^t_{\xi,\ell}:\underline M_\ell\otimes\underline M_\ell
 &\longrightarrow\mathbb Z_\ell,\\
\varphi^t_{\xi,\ell}(x,y)&=u_\ell^{A_K}\bigl(x,\xi_t(y)\bigr).
\end{aligned}
\end{equation}
La forma precedente $\varphi^t_{\xi,\ell}$ es \underline{simétrica}:
\begin{equation}\tag{10.2.5}
\varphi^t_{\xi,\ell}(x,y)=\varphi^t_{\xi,\ell}(y,x).
\end{equation}
En realidad, esta afirmación es válida para todo elemento $\xi$ del
grupo de Néron--Severi de $A_K$, no necesariamente una polarización.
Pero solo cuando $\xi$ es no degenerada, es decir, define una
\underline{isogenia} $A_K\longrightarrow A'_K$, y por tanto también
una isogenia (10.2.3), la relación de simetría (10.2.5) equivale a
(10.2.2). Cuando $\xi$ es una polarización, reforzaremos
considerablemente (10.2.5) en 10.4(b) más abajo.

Señalemos finalmente una tercera interpretación de la relación de
simetría. Consideremos una biextensión (es decir, una correspondencia
divisorial) $\xi$ sobre un par $(A_K,\overline A_K)$ de esquemas
abelianos con reducción semiestable sobre $K$. Define un homomorfismo
\begin{equation*}
A_K\longrightarrow\overline A'_K.
\end{equation*}

%%%%%%%%%%  scan idx 452  |  folio 441  %%%%%%%%%%

Al pasar a los modelos de Néron, este define un homomorfismo sobre los
toros maximales de las fibras especiales y, por tanto, en sentido
inverso, un homomorfismo sobre sus grupos de caracteres:
\begin{equation}\tag{10.2.6}
\xi_t:\overline{\underline M}\longrightarrow\underline M',
\end{equation}
donde $\overline{\underline M}$ se define a partir de
$\overline A_K$ como $\underline M$ se define a partir de $A_K$.
Podemos, por tanto, formar el acoplamiento
\begin{equation}\tag{10.2.7}
\begin{aligned}
\varphi^t_\xi:\underline M_\ell\otimes
 \overline{\underline M}_\ell&\longrightarrow\mathbb Z_\ell,\\
\varphi^t_\xi(x,y)&=u_\ell^{A_K}\bigl(x,\xi_t(y)\bigr).
\end{aligned}
\end{equation}
Observemos de paso que, para un esquema abeliano variable $A_K$ con
reducción semiestable, el grupo constante retorcido $\underline M$ es
un funtor covariante de $A_K$. La formación de los acoplamientos
(10.2.7) asociados a biextensiones de esquemas abelianos con reducción
semiestable sobre $(K,S)$ es entonces compatible, en el sentido
evidente, con la formación de imágenes inversas de biextensiones.

Cuando $\xi$ es la biextensión canónica (``de Weil'') de
$(A_K,A'_K)$, el acoplamiento (10.2.7) es, naturalmente, precisamente
el acoplamiento de monodromía \sisi{(9.2.1)}{(9.1.2)}. Consideremos
ahora la transpuesta ${}^s\xi$ de $\xi$, una biextensión de
$(\overline A_K,A_K)$, y por tanto el acoplamiento
\begin{equation*}
\varphi^t_{{}^s\xi}:\overline{\underline M}_\ell\otimes
 \underline M_\ell\longrightarrow\mathbb Z_\ell.
\end{equation*}
La relación de simetría es
\begin{equation}\tag{10.2.8}
\varphi^t_{{}^s\xi}=\varphi^t_\xi\circ s,
\qquad
s:\overline{\underline M}_\ell\otimes\underline M_\ell
 \xrightarrow{\ \sim\ }
 \underline M_\ell\otimes\overline{\underline M}_\ell,
\end{equation}
donde $s$ es de nuevo el isomorfismo de simetría. Compárese esta
relación con la relación de antisimetría de VIII 2.2.11.
